\documentclass[11pt,a4paper]{article}

\usepackage[a4paper,margin=22mm]{geometry}
\usepackage{fontspec}
\usepackage{babel}
\babelprovide[import, main]{russian}
\babelfont{rm}{Noto Serif}
\babelfont{sf}{Noto Sans}
\usepackage{microtype}
\sloppy
\setlength{\emergencystretch}{4em}

\usepackage{amsmath,amssymb}
\usepackage{booktabs,tabularx}
\usepackage{xcolor}
\usepackage{tikz}
\usetikzlibrary{arrows.meta,positioning,shapes.geometric}
\usepackage{hyperref}

\definecolor{primary}{RGB}{22,62,118}
\definecolor{accent}{RGB}{19,122,99}
\definecolor{warn}{RGB}{176,58,46}
\definecolor{softbg}{RGB}{241,245,251}
\definecolor{warmbg}{RGB}{252,246,238}

\hypersetup{
  colorlinks=true,
  linkcolor=primary,
  citecolor=primary,
  urlcolor=primary
}

\setlength{\parskip}{6.5pt plus 2pt minus 1pt}
\setlength{\parindent}{0pt}

\newcommand{\Pcl}{\mathsf{P}}
\newcommand{\NP}{\mathsf{NP}}
\newcommand{\UP}{\mathsf{UP}}
\newcommand{\PH}{\mathsf{PH}}
\newcommand{\RP}{\mathsf{RP}}
\newcommand{\Modp}{\mathrm{Mod}_p\mathsf{P}}
\newcommand{\SigTwo}{\Sigma_2^{\mathsf{P}}}
\newcommand{\OR}{\mathrm{OR}}
\newcommand{\F}{\mathbb{F}}
\newcommand{\Z}{\mathbb{Z}}
\newcommand{\GI}{GI_1}

\newcommand{\callout}[2]{%
  \vspace{5pt}\noindent\colorbox{softbg}{%
    \parbox{\dimexpr\linewidth-2\fboxsep\relax}{%
      \vspace{4pt}\textbf{\color{primary}#1.}\hspace{0.45em}#2\vspace{4pt}%
    }%
  }\vspace{5pt}%
}

\newcommand{\honestbox}[2]{%
  \vspace{5pt}\noindent\colorbox{warmbg}{%
    \parbox{\dimexpr\linewidth-2\fboxsep\relax}{%
      \vspace{4pt}\textbf{\color{warn}#1.}\hspace{0.45em}#2\vspace{4pt}%
    }%
  }\vspace{5pt}%
}

\title{\textbf{\Large Интерфейс как ресурс познания:\\[4pt]
\large Почему устранение неоднозначности не делает сложные задачи простыми}}
\author{Александр Странд (Aleksandr Strand)\\
\small Независимый исследователь --- научно-популярный препринт по мотивам работы\\
\small\emph{Recursive Switching and Modular Oracle Separations under $\NP = \UP$}}
\date{26 ноября 2026 г.}

\begin{document}
\maketitle

\begin{abstract}
Представьте две вычислительные системы, перед которыми лежит \emph{одна и та же} таблица данных. Ни одна не получает новых скрытых фактов. Внутри мира устранена всякая двусмысленность: у каждого истинного утверждения существует ровно одно уникальное свидетельство. Первая система умеет считать ответы по модулю~$2$,~$3$ или любого другого фиксированного простого числа~$p$, и ей даже разрешено экспоненциально наращивать размер алгебраического поля над этой характеристикой ($\F_3 \to \F_{3^{100}}$). Вторая система умеет менять саму характеристику поля ($\F_3, \F_5, \F_7, \dots$).

В этом тексте без перегрузки формализмом объясняется, что именно мы математически раскопали в препринте \emph{Recursive Switching and Modular Oracle Separations under $\NP = \UP$}: почему первая система \emph{принципиально} не может решить простую на вид задачу о поиске пустой строки, тогда как вторая решает её за полиномиальное время — и почему из этого результата вытекает строгое различие между количественным масштабированием и сменой вычислительного интерфейса. Отдельно и честно проведена граница между тем, что доказано аналитически, что машинно проверено в Lean~4.36.0, а что остаётся открытой проблемой или гипотезой для архитектуры ИИ.
\end{abstract}

\section{Библиотека без двусмысленности}

Чтобы понять суть результата, не нужны сложные определения из теории сложности. Достаточно представить гигантскую библиотеку из $N$ шкафов, в каждом из которых ровно $L$ пронумерованных ячеек (в нашей математической модели $N = L = 2^m$, то есть при $m = 100$ речь идёт о масштабах порядка $10^{30}$).

Библиотека устроена максимально аккуратно: в каждом шкафу лежит \emph{либо ровно один} документ-свидетельство, \emph{либо ни одного}. Никаких дубликатов, никаких противоречивых копий.

Теперь зададим библиотекарю два на вид похожих вопроса:
\begin{enumerate}
  \item \textbf{Вопрос первого типа:} «Лежит ли документ в конкретном шкафу~$u$?»
  \item \textbf{Вопрос второго типа:} «Существует ли во всей библиотеке хотя бы один \emph{полностью пустой} шкаф?»
\end{enumerate}

На первый вопрос легко дать короткое и однозначное доказательство положительного ответа: достаточно открыть нужную ячейку и показать единственный документ. В теории сложности задачи с быстро проверяемым свидетельством образуют класс $\NP$, а их однозначный вариант (где у каждого «да» существует ровно одно принимающее вычисление) — класс $\UP$.

Второй вопрос устроен иначе. Чтобы убедиться, что шкаф~$u$ пуст, нельзя просто показать один документ — нужно исключить \emph{все} $L$ его ячеек:
\[
  \underbrace{\exists u}_{\text{существует шкаф}}
  \quad
  \underbrace{\forall v}_{\text{во всех его ячейках}}
  \quad A(u, v) = 0.
\]
Это чередование кванторов «существует / для всех» выводит задачу на второй уровень полиномиальной иерархии — класс $\SigTwo$.

\begin{figure}[!htbp]
\centering
\begin{tikzpicture}[scale=0.62, every node/.style={font=\sffamily\small}]
  \foreach \r in {0,1,2,3,4} {
    \node[anchor=east, primary] at (-0.3, 4-\r+0.5) {Шкаф \the\numexpr\r+1\relax};
    \foreach \c in {0,1,2,3,4,5,6,7} {
      \draw[gray!50] (\c, 4-\r) rectangle (\c+1, 4-\r+1);
    }
  }
  % Unique green witness cells in rows 1,2,3,5 (row 4 is empty!)
  \fill[accent!75] (2,4) rectangle (3,5);
  \fill[accent!75] (6,3) rectangle (7,4);
  \fill[accent!75] (1,2) rectangle (2,3);
  \fill[accent!75] (5,0) rectangle (6,1);
  % Highlight empty row 4
  \fill[warn!12] (0,1) rectangle (8,2);
  \draw[warn, thick] (0,1) rectangle (8,2);
  \node[anchor=west, warn, font=\sffamily\footnotesize\bfseries] at (8.3, 1.5) {Пустой шкаф ($\forall v:\ A(4,v)=0$)};
  \node[anchor=west, accent!90!black, font=\sffamily\footnotesize] at (8.3, 4.5) {Уникальное свидетельство};
  \draw[->, accent!90!black, thick] (8.2, 4.5) -- (3.1, 4.5);
\end{tikzpicture}
\caption{Модель целевой задачи: в каждом непустом шкафу лежит ровно один документ (зелёная ячейка). Задача $\SigTwo$ — отличить библиотеку без пустых шкафов от библиотеки, где хотя бы один шкаф полностью пуст.}
\end{figure}

Интуитивно кажется, что если мир настолько прозрачен, что \emph{каждый} вопрос из $\NP$ имеет однозначный проверяющий алгоритм ($\NP = \UP$), а вдобавок наблюдателю разрешено мгновенно считать количество принимающих веток по модулю~$2$,~$3$ или любого другого простого числа~$p$, то поиск пустого шкафа должен быстро сломаться. Главный результат нашей работы показывает, что это не так.

\section{Что именно утверждает главная теорема}

В релятивизированной постановке (где все машины обращаются к общему математическому оракулу~$A$, как в классической работе Бейкера, Гилла и Соловея~\cite{BakerGillSolovay1975}) мы строим оракул~$A$ с ациклическим инвариантом уникальных свидетелей $\GI$ (восходящим к конструкции Глассера и Вехзунга~\cite{GlasserWechsung2003}), относительно которого одновременно выполняются два утверждения:
\[
  \NP^A = \UP^A
  \qquad\text{и}\qquad
  \mathcal{L}_A \in \SigTwo{}^{,A} \setminus \bigcup_{p\text{ прост.}} \Modp^A.
\]

Почему здесь важно именно $\SigTwo{}^{,A}$, а не $\NP^A$? При $\NP^A = \UP^A$ каждый язык из $\NP^A$ допускает машину с $0$ или $1$ принимающей ветвью, а числа $0$ и $1$ всегда различны по любому модулю $p \ge 2$. Поэтому $\NP^A$ (и даже весь класс $\Pcl^{\NP^A}$) \emph{автоматически} попадает внутрь $\Modp^A$. Значит, второй уровень $\SigTwo{}^{,A}$ — это самый первый этаж полиномиальной иерархии, на котором вообще возможно сопротивление модульному подсчёту при $\NP^A = \UP^A$.

\begin{table}[!htbp]
\centering
\begin{tabularx}{\linewidth}{lX}
\toprule
\textbf{Ресурс наблюдателя в мире $A$} & \textbf{Судьба задачи «есть ли пустой шкаф» ($\mathcal{L}_A$)} \\
\midrule
Однозначные свидетельства для всего $\NP^A$ ($\NP^A = \UP^A$) & Выполнено по построению оракула \\
Подсчёт по модулю $2$, $3$ или любого фикс.\ простого $p$ & \textbf{Недостаточно} для любой полиномиальной машины \\
Алгебраический доступ над $\F_p \to \F_{p^k}$ (фикс.\ $p$) & \textbf{Недостаточно} в детерминированном режиме \\
Добавление случайности ($\RP^{\widetilde{A}_{\F_3}}$) & \textbf{Разрешимо} (изоляция Валианта--Вазирани) \\
Смена характеристики поля ($\F_3, \F_5, \F_7, \dots$) & \textbf{Разрешимо детерминированно} (через КТО) \\
\bottomrule
\end{tabularx}
\caption{Что именно преодолевает и не преодолевает барьер в построенном мире $A$.}
\end{table}

Более того, речь идёт не об одном экзотическом оракуле-исключении: по теореме Бэра множество таких разделяющих оракулов образует \emph{плотное $G_\delta$} в компактном пространстве всех $\GI$-совместимых оракулов $X_{\GI}$. Какой бы допустимый конечный фрагмент библиотеки вы ни зафиксировали заранее, его всегда можно достроить до мира, в котором разделение сохраняется.

\section{Как это доказано: четыре инженерных механизма}

Главная трудность доказательства состоит в том, что оракул~$A$ не просто хранит базовую матрицу шкафов. Чтобы обеспечить $\NP^A = \UP^A$, оракул обязан содержать \emph{бесконечную рекурсивную башню} верхних слоёв: каждый верхний слой кодирует уникальными маркерами ответы на $\NP$-вопросы о предыдущих слоях. Если разместить маркеры неосторожно (например, всегда в первой ячейке), то уже второй слой башни тривиально вычислит пустоту любого шкафа и разрушит разделение.

Чтобы удержать башню под контролем, в доказательстве работают четыре сцеплённых механизма.

\begin{figure}[!htbp]
\centering
\begin{tikzpicture}[
  node distance=4mm,
  box/.style={draw=primary!70, fill=softbg, rounded corners=3pt, text width=2.95cm, align=center, minimum height=1.4cm, font=\sffamily\footnotesize},
  arr/.style={->, >=Stealth, thick, primary}
]
  \node[box] (s1) {\textbf{1. Singleton Switching}\\Упрощаем базовый слой, не ломая мир без пустых шкафов};
  \node[box, right=of s1] (s2) {\textbf{2. Память и захват}\\Каждый блокирующий раунд открывает \emph{новый} маркер};
  \node[box, right=of s2] (s3) {\textbf{3. Префиксный поиск}\\Рекурсия по $O(\log m)$ слоям без взрыва ширины};
  \node[box, right=of s3] (s4) {\textbf{4. Двойственный $\Lambda_R$}\\$\deg \le (p{-}1)Td_h < R = \deg(\OR_R)$};
  \draw[arr] (s1) -- (s2);
  \draw[arr] (s2) -- (s3);
  \draw[arr] (s3) -- (s4);
\end{tikzpicture}
\caption{Четыре звена доказательной цепочки от базовой матрицы до невозможности модульного решателя.}
\end{figure}

\textbf{Шаг 1. Управляемая деформация (Singleton-Preserving Switching).} Классический метод Хостада~\cite{Hastad1986} (и его варианты над большими алфавитами Хирахары и Нанашимы~\cite{HiraharaNanashima2021}) упрощает сложные формулы, случайно фиксируя переменные в $0$ и $1$. Но нам \emph{нельзя} фиксировать шкафы пустыми, иначе исчезнет сам мир «без пустых шкафов». Поэтому мы используем одностороннюю проекцию: часть шкафов фиксируется непустыми, а оставшиеся остаются свободными. Пустые состояния появляются лишь во вспомогательной мере кодирования для доказательства инъективности, но никогда не фиксируются в итоговом мире.

\textbf{Шаг 2. Постоянная память препятствий (Persistent Marker Knowledge и Trace Capture).} На верхних слоях положительные условия формулы могут работать как «прицел», направляющий отрицательные запросы точно в скрытые маркеры. Чтобы сломать этот прицел, алгоритм запоминает все уже раскрытые истинные маркеры. Тогда каждый новый блокирующий раунд обязан обнаружить \emph{принципиально новый} маркер, не пересекающийся с ранее раскрытыми положительными строками. Случайное разделение строк с вероятностью $[e(w+1)]^{-s}$ захватывает всю такую трассу целиком.

\textbf{Шаг 3. Рекурсия по префиксам без взрыва.} В силу строгой ацикличности $\GI$ полиномиальное вычисление за время $T = \mathrm{poly}(m)$ успевает затронуть не более $h = O(\log m)$ вложенных слоёв башни. Вместо раскрытия верхних формул обратно в гигантские базовые формулы мы ищем первый допустимый терм бинарным поиском по префиксным предикатам предыдущего слоя. В итоге все оракульные биты до длины $T$ одновременно сводятся к деревьям решений квазиполиномиальной глубины $d_h \le 2^{O(\log^4 m)}$ над остаточным свободным кубом размерности $R \ge 2^{m - O(\log^5 m)}$.

\textbf{Шаг 4. Алгебраический потолок степени.} Любая машина класса $\Modp^A$ сводится (по малой теореме Ферма) к целочисленному полиному $C_M(A)^{p-1} \bmod p$ степени не выше $(p-1)Td_h$ от оставшихся $R$ свободных переменных. Но сама функция $\OR_R(b) = 1 - \prod_{i=1}^R (1 - b_i)$ требует полной степени $R$. Знакопеременный функционал
\[
  \Lambda_R(f) = \sum_{b \in \{0,1\}^R} (-1)^{|b|} f(b)
\]
обнуляет \emph{любой} полином степени меньше $R$, но на целевой функции даёт $\Lambda_R(\OR_R) = -1 \not\equiv 0 \pmod p$ для любого $p \ge 2$. Поскольку $(p-1)Td_h = 2^{O(\log^4 m)} \ll 2^{m - O(\log^5 m)} = R$, на остаточном кубе гарантированно существует допустимый мир, в котором машина ошибается.

\section{Пять эпистемологических сдвигов}

Если отойти от технических деталей, в этой работе математически разведены пять понятий, которые в повседневной (и инженерной) интуиции часто склеивают в одно.

\textbf{1. Однозначность знания $\ne$ простота глобального вывода.} Устранение двусмысленности ($\NP^A = \UP^A$) гарантирует, что у каждого локального факта ровно один сертификат. Но это никак не удешевляет глобальный вывод об отсутствии сертификатов.

\textbf{2. Одинаковая информация $\ne$ одинаковая вычислительная мощность.} В модели алгебраизации Ааронсона--Вигдерсона~\cite{AaronsonWigderson2009} (и свежих мультилинейных барьерах Чена, Ху и Жэня~\cite{ChenHuRen2026}) мультилинейное продолжение $\widetilde{A}$ однозначно определяется таблицей $A$ и не добавляет ни одного нового независимого бита. Тем не менее интерфейс доступа становится самостоятельным вычислительным ресурсом.

\begin{figure}[!htbp]
\centering
\begin{tikzpicture}[
  node distance=8mm,
  obj/.style={draw=primary, fill=softbg, rounded corners=4pt, text width=3.4cm, align=center, minimum height=1.2cm, font=\sffamily\small\bfseries},
  obsA/.style={draw=warn, fill=warmbg, rounded corners=4pt, text width=5.4cm, align=center, minimum height=1.5cm, font=\sffamily\footnotesize},
  obsB/.style={draw=accent, fill=accent!10, rounded corners=4pt, text width=5.4cm, align=center, minimum height=1.5cm, font=\sffamily\footnotesize},
  arr/.style={->, >=Stealth, thick}
]
  \node[obj] (A) {Один и тот же оракул $A$\\{\normalfont\footnotesize одинаковая исходная информация}};
  \node[obsA, below left=10mm and -12mm of A] (left) {\textbf{\color{warn}Наблюдатель A (фикс.\ характеристика $p$)}\\Масштабирует поле: $\F_3 \to \F_{3^{10}} \to \F_{3^{100}}$\\Детерминированно $\mathcal{L}_A$ \textbf{недостижима}};
  \node[obsB, below right=10mm and -12mm of A] (right) {\textbf{\color{accent}Наблюдатель B (смена характеристик)}\\Меняет арифметику: $\F_3, \F_5, \F_7, \dots$ (КТО)\\Детерминированно $\mathcal{L}_A$ \textbf{решается в $\Pcl$}};
  \draw[arr, warn] (A.south west) -- (left.north);
  \draw[arr, accent] (A.south east) -- (right.north);
\end{tikzpicture}
\caption{Разделение количественного масштабирования (рост размера поля при фиксированной характеристике) и структурного преобразования интерфейса (смена характеристики).}
\end{figure}

\textbf{3. Масштабирование внутри языка $\ne$ смена языка.} Наращивание размера поля $\F_p \to \F_{p^k}$ в детерминированном режиме не даёт никакого прироста мощности против $\mathcal{L}_A$ (поскольку $\Pcl^{\widetilde{A}^{[p]}} \subseteq \Modp^A$). Барьер снимает только качественное изменение: либо добавление случайности, либо разрешение менять характеристику поля.

\textbf{4. Локальная простота $\ne$ глобальная сертифицируемость.} Каждый отдельный наблюдаемый бит после проекции имеет малую степень $d_h$, но их полиномиальные комбинации всё равно не дотягивают до степени $R$ глобального свойства $\OR_R$.

\textbf{5. Аномалия $\ne$ типичность.} Плотное $G_\delta$ в $X_{\GI}$ показывает, что разделение — не хрупкий искусственный контрпример, а топологически типичное поведение допустимых миров (хотя топологическую типичность не следует путать с вероятностной мерой~$1$).

\section{Что это значит (и чего не значит) для архитектуры ИИ}

Результат оракульной теории сложности сам по себе \emph{не доказывает} теорем о нейросетях или LLM. Но он даёт строгую модельную аналогию для инженерной интуиции, с которой сталкиваются разработчики агентных систем.

\callout{Иллюстрация с 6 пустыми шкафами}{Представьте 8 шкафов, из которых 6 пусты. Наблюдатель с интерфейсом $\bmod 2$ видит $6 \bmod 2 = 0$; наблюдатель с интерфейсом $\bmod 3$ видит $6 \bmod 3 = 0$. Даже десять параллельных агентов с тем же интерфейсом $\bmod 2$ не устранят неоднозначность простым увеличением бюджета токенов. Но смена интерфейса на $\bmod 5$ мгновенно даёт $6 \bmod 5 = 1 \ne 0$.}

В терминах архитектуры ИИ-агентов это разводит два режима работы:
\begin{itemize}
  \item \textbf{Количественное масштабирование:} больше токенов, попыток, памяти и копий агента внутри фиксированного набора операций.
  \item \textbf{Структурное масштабирование:} обнаружение ограничений текущего представления и проверяемый переход к другому вычислительному интерфейсу при сохранении критерия истинности исходной задачи.
\end{itemize}

Для кумулятивного исследовательского агента отсюда возникает естественное деление памяти на четыре слоя: \emph{фактическую} (что узнали), \emph{процедурную} (как сделали), \emph{эпистемическую} (почему данный класс методов упирается в барьер) и \emph{структурную} (какая смена представления сняла барьер).

\honestbox{Честная оговорка для ИИ}{В математической модели смена характеристики даётся как готовый оракульный ресурс. В реальной системе ИИ новый интерфейс приходится строить и проверять, а стоимость его построения обязана входить в общий бюджет. Поэтому перенос этого принципа в архитектуру агентов — мотивированная исследовательская гипотеза, требующая отдельного контролируемого эксперимента, а не доказанное превосходство.}

\section{Честная карта границ: чем эта работа не является}

В научно-популярных изложениях математических результатов особенно важно не размывать статус доказанного. Ниже — точный реестр того, что установлено, а что нет.

\begin{table}[!htbp]
\centering
\begin{tabularx}{\linewidth}{Xl}
\toprule
\textbf{Утверждение} & \textbf{Фактический статус} \\
\midrule
Решение нерелятивизированной проблемы $\Pcl$ vs.\ $\NP$~\cite{Cook2000} & \textbf{Нет} \\
Доказательство бесконечности $\PH^A$ при $\NP^A = \UP^A$ (Grochow~\cite{Grochow2023}) & \textbf{Нет (открыто)} \\
Универсальный обход барьера алгебраизации & \textbf{Нет} \\
Доказательство превосходства новой архитектуры LLM & \textbf{Нет (гипотеза)} \\
Полная сквозная Lean-формализация бесконечного оракула & \textbf{Пока нет} \\
Аналитическое доказательство $\SigTwo{}^{,A} \not\subseteq \bigcup_p \Modp^A$ при $\NP^A = \UP^A$ & \textbf{Представлено} \\
Машинная проверка выбранных ядер и конечного моста в Lean~4.36.0 & \textbf{7 модулей, 45/45 теорем} \\
\bottomrule
\end{tabularx}
\caption{Разграничение доказанных результатов, частичной формализации и открытых вопросов.}
\end{table}

Отдельно подчеркнём границу формальной верификации: семь 0-import модулей Lean~4.36.0 (45 аудированных теорем, 0 ошибок, 0 \texttt{sorryAx}, 0 \texttt{Classical.choice}) машинно проверяют выбранные комбинаторные и алгебраические ядра, а также \emph{условный} конечный мост (\texttt{UniformRecursiveBridge.lean}), но полная вероятностная рекурсия и построение бесконечного оракула остаются аналитическими, а внешнее научное рецензирование ещё предстоит.

\section{Вместо заключения}

Обычный уровень решения трудной задачи — искать ответ внутри заданного языка. Следующий уровень — понять, почему сам выбранный язык не способен отличить нужные состояния. И только на третьем уровне появляется возможность задать главный конструктивный вопрос: \emph{какое минимальное изменение интерфейса делает прежде недоступное знание вычислимым?}

Именно это разделение — между объёмом данных и геометрией доступа к ним — и составляет главный смысл того, что нам удалось раскопать.

\begin{thebibliography}{9}

\bibitem{AaronsonWigderson2009}
S.~Aaronson and A.~Wigderson.
Algebrization: A new barrier in complexity theory.
\emph{ACM Transactions on Computation Theory}, 1(1):2:1--2:54, 2009.
\href{https://www.scottaaronson.com/papers/alg.pdf}{PDF}.

\bibitem{BakerGillSolovay1975}
T.~Baker, J.~Gill, and R.~Solovay.
Relativizations of the $\Pcl = ?\NP$ question.
\emph{SIAM Journal on Computing}, 4(4):431--442, 1975.
\href{https://epubs.siam.org/doi/10.1137/0204037}{DOI}.

\bibitem{ChenHuRen2026}
L.~Chen, Y.~Hu, and H.~Ren.
New algebrization barriers to circuit lower bounds via communication complexity of Missing-String.
In \emph{ITCS 2026}, 2026.
\href{https://arxiv.org/abs/2511.14038}{arXiv:2511.14038}.

\bibitem{Cook2000}
S.~Cook.
The $\Pcl$ versus $\NP$ problem.
Clay Mathematics Institute, 2000.
\href{https://www.claymath.org/wp-content/uploads/2022/06/pvsnp.pdf}{PDF}.

\bibitem{GlasserWechsung2003}
C.~Gla\ss er and G.~Wechsung.
Relativizing function classes.
\emph{Journal of Universal Computer Science}, 9(1):34--50, 2003.
\href{https://lib.jucs.org/issue/4069/}{JUCS 9(1)}.

\bibitem{Grochow2023}
J.~A.~Grochow.
Polynomial-time axioms of choice and polynomial-time cardinality.
\emph{Theory of Computing Systems}, 67, 2023.
\href{https://arxiv.org/abs/2301.07123}{arXiv:2301.07123}.

\bibitem{Hastad1986}
J.~H\aa stad.
Almost optimal lower bounds for small depth circuits.
In \emph{STOC 1986}, pp.~6--20, 1986.

\bibitem{HiraharaNanashima2021}
S.~Hirahara and M.~Nanashima.
On worst-case learning in relativized Heuristica.
In \emph{FOCS 2021}, pp.~751--758, 2021.
\href{https://eccc.weizmann.ac.il/report/2021/161/download}{ECCC TR21-161}.

\end{thebibliography}

\end{document}
